\section{Effective descent: restriction of the glued sheaf to a chart}
\label{mk-chap:Picard_GlueDescent}

\providecommand{\ShMod}{\mathrm{SheafOfModules}}

The gluing primitive \ref{mk-def:scheme_modules_glue} produces, out of a glue datum
\(D\) of schemes (index set \(J\), charts \(U_i\), overlaps \(V_{ij}\), overlap
immersions \(f_{ij} : V_{ij} \hookrightarrow U_i\), transition isomorphisms
\(t_{ij} : V_{ij} \xrightarrow{\sim} V_{ji}\), glued scheme \(X\) with chart
immersions \(\iota_i : U_i \hookrightarrow X\)), a sheaf of \(\mathcal{O}_X\)-modules
\(\mathcal{M}\) as the descent equalizer
\[
  \mathcal{M} \;=\; \operatorname{eq}(a, b)
    \;\hookrightarrow\; \prod_{i} (\iota_i)_{*}\,\mathcal{M}_i
\]
of two parallel maps into the overlap product. That construction makes
\(\mathcal{M}\) but does \emph{not} yet consume the cocycle conditions (C1)/(C2) on
the transition family \((g_{ij})\). This chapter discharges the remaining content of
descent: it shows that restricting the glued sheaf back to a chart recovers the
input sheaf there. Concretely, the canonical \emph{restriction morphism}
\(\iota_i^{*}\mathcal{M} \to \mathcal{M}_i\) -- the adjoint transpose along the chart
immersion \(\iota_i\) of the \(i\)-th equalizer projection -- is an isomorphism, and
it is here that (C1)/(C2) are used. The engine is that pullback along an open
immersion is, up to natural isomorphism, the site-level restriction functor
(\ref{mk-lem:modules_restrictFunctorIsoPullback_mathlib}); being a right adjoint it
preserves the descent equalizer and the pushforward product, so the chart
restriction of \(\mathcal{M}\) is computed factor by factor. The non-formal input is
the \emph{overlap base change} \(\beta_{ij}\), which identifies the restricted
pushforward factor \(\iota_i^{*}(\iota_j)_{*}\mathcal{M}_j\) with a pushforward of a
restriction of \(\mathcal{M}_j\) to the overlap, using the cartesianness of the
chart--overlap square.

\section{The descent equalizer and its overlap base change}
\label{mk-sec:gd_equalizer}

\begin{definition}[The pushforward product carrying the descent equalizer]\label{mk-def:gr_glue_equalizer}
  \dcref{custom}

  \uses{mk-def:scheme_modules_glue}
  Let \(D\) be a glue datum and \((\mathcal{M}_i)_{i \in J}\) a family of sheaves of
  modules with \(\mathcal{M}_i\) on the chart \(U_i\). Write
  \[
    P \;=\; \prod_{i} (\iota_i)_{*}\,\mathcal{M}_i
    \qquad\text{and}\qquad
    Q \;=\; \prod_{(i,j)} (f_{ij} \circ \iota_i)_{*}\,\bigl(f_{ij}^{*}\mathcal{M}_i\bigr)
  \]
  for the \emph{pushforward product} and the \emph{overlap product} of sheaves of
  \(\mathcal{O}_X\)-modules. The two descent legs of \ref{mk-def:scheme_modules_glue},
  \[
    P \;\underset{b}{\overset{a}{\rightrightarrows}}\; Q,
  \]
  are: on the \((i,j)\)-component, leg \(a\) restricts the \(i\)-th chart factor to
  the overlap \(V_{ij}\) using the unit of the pullback--pushforward adjunction along
  \(f_{ij}\) and the pushforward-composition comparison
  \((\iota_i)_{*} \cong (f_{ij} \circ \iota_i)_{*} \circ f_{ij}^{*}\); leg \(b\)
  restricts the \(j\)-th chart factor, transports it across the inverse transition
  isomorphism \(g_{ij}^{-1}\), and reindexes the immersion via the glue condition
  \((t_{ij} \circ f_{ji}) \circ \iota_j = f_{ij} \circ \iota_i\). The glued sheaf
  \(\mathcal{M} = \operatorname{eq}(a,b)\) is the equalizer of this parallel pair;
  \(P\) is the object it embeds into and \(Q\) the object detecting the overlap
  agreement. The product \(P\) is the carrier on which the \(i\)-th projection (and,
  through it, the restriction morphism) is defined.
\end{definition}

\begin{lemma}[Overlap base change of the chart square (\(\beta_{ij}\))]\label{mk-lem:glueOverlapBaseChangeIso}
  \dcref{custom}

  \uses{mk-lem:modules_restrictFunctorIsoPullback_mathlib}
  For chart indices \(i, j \in J\) there is a natural isomorphism of functors
  \(\ShMod(\mathcal{O}_{U_j}) \to \ShMod(\mathcal{O}_{U_i})\),
  \[
    \beta_{ij} \;:\;
      \iota_i^{*}\,(\iota_j)_{*}
        \;\xrightarrow{\ \sim\ }\;
      (f_{ij})_{*}\,(t_{ij} \circ f_{ji})^{*},
  \]
  identifying, for any sheaf \(\mathcal{N}\) on the chart \(U_j\), the restriction to
  the chart \(U_i\) of the pushforward \((\iota_j)_{*}\mathcal{N}\) with the
  pushforward along \(f_{ij}\) of the restriction of \(\mathcal{N}\) to the overlap
  \(V_{ij}\) along \(t_{ij} \circ f_{ji}\). Here the chart restriction \(\iota_i^{*}\)
  and the overlap restriction \((t_{ij} \circ f_{ji})^{*}\) are realised as the
  site-level restriction functors, isomorphic to the geometric pullbacks by
  \ref{mk-lem:modules_restrictFunctorIsoPullback_mathlib}.
\end{lemma}
\begin{proof}
  \uses{mk-lem:modules_restrictFunctorIsoPullback_mathlib, mk-lem:gr_overlap_appIso_compat}
    The isomorphism is the cartesianness of the chart--overlap square, made site-level.
  The square
  \[
    \begin{array}{ccc}
      V_{ij} & \xrightarrow{\ t_{ij} \circ f_{ji}\ } & U_j \\
      \ \downarrow{\scriptstyle f_{ij}} & & \downarrow{\scriptstyle \iota_j} \\
      U_i & \xrightarrow{\ \ \iota_i\ \ } & X
    \end{array}
  \]
  is cartesian, and its cartesianness is visible already on underlying opens: for an
  open \(V \subseteq U_i\),
  \[
    \iota_j^{-1}\bigl(\iota_i(V)\bigr)
      \;=\; (t_{ij} \circ f_{ji})\bigl(f_{ij}^{-1}(V)\bigr),
  \]
  a point of \(U_j\) lands in the image of \(V\) under \(\iota_i\) precisely when it
  comes, via the glue relation, from a point of the overlap lying over \(V\). Hence
  the two composite opens functors
  \[
    (\iota_i)_{!} \circ \iota_j^{-1}
      \;=\;
    f_{ij}^{-1} \circ (t_{ij} \circ f_{ji})_{!}
  \]
  agree on objects, and -- morphisms of the opens poset being proof-irrelevant -- as
  functors. Both \(\iota_i^{*}(\iota_j)_{*}\) and \((f_{ij})_{*}(t_{ij}\circ
  f_{ji})^{*}\) are site-level pushforwards along these two (equal) opens functors, so
  the comparison is assembled from the pushforward-composition isomorphisms together
  with the natural isomorphism transporting one opens functor to the other:
  \[
    \beta_{ij}
      \;=\;
    (\text{pushforwardComp})
      \;\circ\;
    (\text{pushforward of the opens-functor equality})
      \;\circ\;
    (\text{pushforwardCongr})
      \;\circ\;
    (\text{pushforwardComp})^{-1}.
  \]
  The one substantive coherence to check is that the comparison glued from the
  object-level opens equality is natural in the open \(V\): on the section level the
  two restriction maps induced by the equal composite functors coincide because the
  connecting morphisms are the unique arrows of the opens poset between the identified
  opens. This is the pointwise compatibility that completes the \(\mathrm{pushforwardCongr}\)
  factor.
\end{proof}

\section{The restriction morphism and effective descent}
\label{mk-sec:gd_descent}

\begin{lemma}[The chart restriction morphism of the glued sheaf]\label{mk-lem:glueRestrictionHom}
  \dcref{custom}

  \uses{mk-def:gr_glue_equalizer, mk-def:scheme_modules_glue}
  Let \(\mathcal{M}\) be the glued sheaf of \ref{mk-def:scheme_modules_glue}. The
  \(i\)-th \emph{projection} of \(\mathcal{M}\),
  \[
    \pi_i \;:\; \mathcal{M} \longrightarrow (\iota_i)_{*}\,\mathcal{M}_i,
  \]
  is the equalizer inclusion \(\mathcal{M} \hookrightarrow P\)
  (\ref{mk-def:gr_glue_equalizer}) followed by the \(i\)-th product projection. Its
  adjoint transpose along the chart immersion \(\iota_i\), under the
  pullback--pushforward adjunction \(\iota_i^{*} \dashv (\iota_i)_{*}\), is the
  \emph{restriction morphism}
  \[
    \mathrm{r}_i \;:\; \iota_i^{*}\,\mathcal{M} \longrightarrow \mathcal{M}_i .
  \]
  Equivalently, \(\mathrm{r}_i\) is the unique morphism whose pushforward, precomposed
  with the unit, recovers \(\pi_i\).
\end{lemma}
\begin{proof}
  \uses{mk-def:gr_glue_equalizer, mk-def:scheme_modules_glue}
  The projection \(\pi_i\) is well-defined by \ref{mk-def:gr_glue_equalizer}: the glued
  sheaf is the equalizer \(\operatorname{eq}(a,b) \hookrightarrow P\), and composing
  the inclusion with the \(i\)-th projection of the product \(P = \prod_i
  (\iota_i)_{*}\mathcal{M}_i\) gives a map to \((\iota_i)_{*}\mathcal{M}_i\). The
  restriction morphism is then its image under the natural bijection
  \[
    \operatorname{Hom}_{\mathcal{O}_X}\bigl(\mathcal{M}, (\iota_i)_{*}\mathcal{M}_i\bigr)
      \;\cong\;
    \operatorname{Hom}_{\mathcal{O}_{U_i}}\bigl(\iota_i^{*}\mathcal{M}, \mathcal{M}_i\bigr)
  \]
  supplied by the adjunction \(\iota_i^{*} \dashv (\iota_i)_{*}\). The stated
  equivalent description is the unit--counit form of this transpose.
\end{proof}

\begin{theorem}[Effective descent: the chart restriction morphism is an isomorphism]\label{mk-thm:isIso_glueRestrictionHom}
  \dcref{custom}

  \uses{mk-lem:glueRestrictionHom, mk-lem:glueOverlapBaseChangeIso,
    mk-lem:isLimitPullbackCone_mathlib, mk-lem:eq_of_locally_eq_mathlib,
    mk-def:gr_glueRestrictionInv, mk-lem:gr_glueRestrict_hom_ext,
    mk-lem:gr_glueChartComponent_self_counit,
    mk-lem:gr_glueRestrictionHom_glueChartComponent,
    mk-lem:gr_glueRestrictionInv_proj}
  Assume the transition family \((g_{ij})\) satisfies the cocycle conditions (C1) and
  (C2) of \ref{mk-def:scheme_modules_glue}. Then for every chart index \(i\) the
  restriction morphism \(\mathrm{r}_i : \iota_i^{*}\mathcal{M} \to \mathcal{M}_i\)
  (\ref{mk-lem:glueRestrictionHom}) is an isomorphism of sheaves of
  \(\mathcal{O}_{U_i}\)-modules.
\end{theorem}
\begin{proof}
  \uses{mk-lem:glueRestrictionHom, mk-lem:glueOverlapBaseChangeIso,
    mk-lem:isLimitPullbackCone_mathlib, mk-lem:eq_of_locally_eq_mathlib,
    mk-def:gr_glueRestrictionInv, mk-lem:gr_glueRestrict_hom_ext,
    mk-lem:gr_glueChartComponent_self_counit,
    mk-lem:gr_glueRestrictionHom_glueChartComponent,
    mk-lem:gr_glueRestrictionInv_proj}
    The chart pullback \(\iota_i^{*}\) is naturally isomorphic to the site-level
  restriction functor (\ref{mk-lem:modules_restrictFunctorIsoPullback_mathlib}), which
  is a right adjoint; consequently \(\iota_i^{*}\) \emph{preserves limits}. Applying
  it to the descent equalizer \(\mathcal{M} = \operatorname{eq}(a,b)\) exhibits
  \(\iota_i^{*}\mathcal{M}\) as the equalizer of the restricted legs
  \(\iota_i^{*}a, \iota_i^{*}b\), and applying it to the product
  \(P = \prod_j (\iota_j)_{*}\mathcal{M}_j\) exhibits \(\iota_i^{*}P\) as the product
  \(\prod_j \iota_i^{*}(\iota_j)_{*}\mathcal{M}_j\). The overlap base change
  \(\beta_{ij}\) of \ref{mk-lem:glueOverlapBaseChangeIso} then identifies the \(j\)-th
  factor \(\iota_i^{*}(\iota_j)_{*}\mathcal{M}_j\) with
  \((f_{ij})_{*}\bigl((t_{ij} \circ f_{ji})^{*}\mathcal{M}_j\bigr)\). Under these
  identifications the restricted descent data becomes the data attached to the chart
  \(U_i\) by the cover \(\{f_{ij}\}\) of \(U_i\), and the \(i\)-th factor (where
  \(j = i\)) is, by (C1), \(\mathcal{M}_i\) itself.

  We produce an inverse \(\mathrm{s}_i : \mathcal{M}_i \to \iota_i^{*}\mathcal{M}\).
  Because \(\iota_i^{*}\mathcal{M}\) is the equalizer of the restricted legs, it
  suffices to give a morphism \(\mathcal{M}_i \to \iota_i^{*}P\) -- equivalently a
  compatible family of morphisms into the factors -- that equalizes
  \(\iota_i^{*}a\) and \(\iota_i^{*}b\). The \(j\)-th component is
  \[
    \mathcal{M}_i
      \xrightarrow{\ \eta_{f_{ij}}\ }
    (f_{ij})_{*}\,f_{ij}^{*}\mathcal{M}_i
      \xrightarrow{\ (f_{ij})_{*} g_{ij}\ }
    (f_{ij})_{*}\,(t_{ij} \circ f_{ji})^{*}\mathcal{M}_j
      \xrightarrow{\ \beta_{ij}^{-1}\ }
    \iota_i^{*}(\iota_j)_{*}\mathcal{M}_j ,
  \]
  the unit of the pullback--pushforward adjunction along \(f_{ij}\), followed by the
  pushforward of the transition isomorphism \(g_{ij}\), followed by the inverse of the
  overlap base change. That this family equalizes the two restricted legs is exactly
  the triple-overlap multiplicativity (C2), transported through \(\beta\); and the
  composite \(\mathrm{r}_i \circ \mathrm{s}_i = \mathrm{id}_{\mathcal{M}_i}\) reduces,
  via the \(j = i\) component, to the self-identity (C1) together with the counit
  triangle identity of the adjunction. The reverse composite
  \(\mathrm{s}_i \circ \mathrm{r}_i = \mathrm{id}\) holds because both sides are
  equalizer maps with the same components into \(P\); two such agree once their
  factors agree, which is the separatedness of the sheaf equalizer over the chart
  cover. The limit computation underlying these equalities is the sheaf
  Mayer--Vietoris pullback square (\ref{mk-lem:isLimitPullbackCone_mathlib}), and the
  final identification of sections is the uniqueness-over-a-cover principle
  (\ref{mk-lem:eq_of_locally_eq_mathlib}). Hence \(\mathrm{r}_i\) is an isomorphism.
\end{proof}

\section{Construction of the inverse: the descent obligations}
\label{mk-sec:gd_inverse}


\subsection{Section-level compatibilities of the overlap square}

\begin{lemma}[Transport of \(\mathrm{appLE}\) along equal morphisms]\label{mk-lem:gr_appLE_congr_mor}
  \dcref{custom}

  Let \(f, g : A \to B\) be morphisms of schemes with \(f = g\), let \(U\) be an open of
  \(B\) and \(W\) an open of \(A\) with \(W \subseteq f^{-1}(U)\) (equivalently
  \(W \subseteq g^{-1}(U)\)). Then the two induced section maps
  \(\Gamma(B, U) \to \Gamma(A, W)\) along \(f\) and along \(g\) coincide:
  \(f^{\mathrm{LE}}_{U,W} = g^{\mathrm{LE}}_{U,W}\).
\end{lemma}
\begin{proof}

Substituting the equality \(f = g\) reduces both sides to the same map; the two
  open-inclusion witnesses \(W \subseteq f^{-1}(U)\) and \(W \subseteq g^{-1}(U)\) are
  proof-irrelevant.
\end{proof}

\begin{lemma}[Structure-sheaf compatibility of the chart--overlap square]\label{mk-lem:gr_overlap_appIso_compat}
  \dcref{custom}

  \uses{mk-lem:gr_appLE_congr_mor, mk-def:scheme_modules_glue}
  For chart indices \(i, j\) and an open \(V \subseteq U_i\), the two composite
  section maps
  \[
    \Gamma(U_i, V)\;\longrightarrow\;
      \Gamma\bigl(U_j,\ (t_{ij}\circ f_{ji})\bigl(f_{ij}^{-1}(V)\bigr)\bigr)
  \]
  of the chart--overlap square agree. The first goes ``through the glued scheme'':
  the inverse of the comparison isomorphism of \(\iota_i\) on \(V\), the section map of
  \(\iota_j\) over \(\iota_i(V)\), and the restriction along the opens identity
  \(\iota_j^{-1}(\iota_i(V)) = (t_{ij}\circ f_{ji})(f_{ij}^{-1}(V))\). The second goes
  ``through the overlap'': the section map of \(f_{ij}\) over \(V\) followed by the
  inverse comparison isomorphism of \(t_{ij}\circ f_{ji}\).
\end{lemma}
\begin{proof}
  \uses{mk-lem:gr_appLE_congr_mor, mk-def:scheme_modules_glue}
  Both composites are \(\mathrm{appLE}\)s of the two composite morphisms of the square,
  which are equal by the glue condition \((t_{ij}\circ f_{ji})\circ \iota_j =
  f_{ij}\circ \iota_i\) of \ref{mk-def:scheme_modules_glue}. Moving the two comparison
  isomorphisms to \(\mathrm{appLE}\) form (using \(\mathrm{appIso}\) as an
  \(\mathrm{appLE}\), the composition law \(g^{\mathrm{LE}}\circ f^{\mathrm{LE}} =
  (f\circ g)^{\mathrm{LE}}\), and the absorption of presheaf restrictions into the
  open-inclusion witness) reduces the claim to the equality of the two
  \(\mathrm{appLE}\)s of the equal composites, which is \ref{mk-lem:gr_appLE_congr_mor}.
\end{proof}

\begin{lemma}[Sections of the overlap base change, inverse side]\label{mk-lem:gr_overlapBaseChange_inv_app}
  \dcref{custom}

  \uses{mk-lem:glueOverlapBaseChangeIso}
  For a sheaf \(\mathcal{N}\) on \(U_j\) and an open \(V \subseteq U_i\), the inverse
  of the overlap base change \(\beta_{ij}^{-1}\), evaluated on sections over \(V\), is
  the restriction map of \(\mathcal{N}\) along the opens identity
  \(\iota_j^{-1}(\iota_i(V)) = (t_{ij}\circ f_{ji})(f_{ij}^{-1}(V))\) of
  \ref{mk-lem:glueOverlapBaseChangeIso}:
  \[
    (\beta_{ij}^{-1})_{\mathcal{N}}\big|_V
      \;=\;
    \mathcal{N}\!\left(\,(t_{ij}\circ f_{ji})(f_{ij}^{-1}V) \;=\; \iota_j^{-1}(\iota_i V)\,\right).
  \]
\end{lemma}
\begin{proof}
  \uses{mk-lem:glueOverlapBaseChangeIso}
  Each of the four functorial factors of \(\beta_{ij}\) (\ref{mk-lem:glueOverlapBaseChangeIso})
  acts on sections over \(V\) either as the identity or as a single presheaf
  restriction along the opens identity; composing them leaves exactly the named
  restriction map. The verification is pointwise on sections.
\end{proof}

\begin{lemma}[Sections of the overlap base change, forward side]\label{mk-lem:gr_overlapBaseChange_hom_app}
  \dcref{custom}

  \uses{mk-lem:glueOverlapBaseChangeIso}
  Symmetrically, the forward overlap base change \(\beta_{ij}\), evaluated on sections
  over an open \(V \subseteq U_i\), is the restriction map of \(\mathcal{N}\) along the
  reverse opens identity \(\iota_j^{-1}(\iota_i(V)) = (t_{ij}\circ f_{ji})(f_{ij}^{-1}(V))\).
\end{lemma}
\begin{proof}
  \uses{mk-lem:glueOverlapBaseChangeIso}
  Identical to \ref{mk-lem:gr_overlapBaseChange_inv_app} with the roles of source and
  target opens exchanged: pointwise on sections, the composite of the four factors is
  the single presheaf restriction along the (reversed) opens identity.
\end{proof}

\subsection{Bridges between the geometric and the site-level adjunction}

The chart pullback \(\iota^{*}\) is built by sheafification and carries the
\emph{geometric} adjunction \(\iota^{*} \dashv \iota_{*}\); along an open immersion the
\emph{site-level} restriction functor \(\iota^{\mathrm{res}}\) carries a second
adjunction with concrete unit and counit. They share the right adjoint \(\iota_{*}\),
and \ref{mk-lem:modules_restrictFunctorIsoPullback_mathlib} is the uniqueness comparison
of left adjoints between them. The next blocks transport units, counits and
congruence casts across that comparison, so that the descent obligations can be
discharged by section-level computations.

\begin{lemma}[Unit comparison across the adjunctions]\label{mk-lem:gr_restrictAdjunction_unit_iso}
  \dcref{custom}

  \uses{mk-lem:modules_restrictFunctorIsoPullback_mathlib}
  For an open immersion \(f\) and a sheaf \(\mathcal{N}\) on the target, the geometric
  unit is the site-level unit followed by the pushforward of the left-adjoint
  comparison isomorphism:
  \[
    \eta^{\mathrm{res}}_{\mathcal{N}}\ \circ\
      f_{*}\bigl((\text{comparison})_{\mathcal{N}}\bigr)
    \;=\; \eta^{\mathrm{geo}}_{\mathcal{N}}.
  \]
\end{lemma}
\begin{proof}
  \uses{mk-lem:modules_restrictFunctorIsoPullback_mathlib}
  This is the general identity relating the unit of an adjunction to the unit of any
  other adjunction with the same right adjoint, transported by the uniqueness
  isomorphism of left adjoints (\ref{mk-lem:modules_restrictFunctorIsoPullback_mathlib}).
\end{proof}

\begin{lemma}[Counit comparison across the adjunctions]\label{mk-lem:gr_restrictIsoPullback_counit}
  \dcref{custom}

  \uses{mk-lem:modules_restrictFunctorIsoPullback_mathlib}
  Dually, the left-adjoint comparison isomorphism on \(f_{*}\mathcal{N}\) followed by
  the geometric counit equals the site-level counit:
  \[
    (\text{comparison})_{f_{*}\mathcal{N}}\ \circ\ \varepsilon^{\mathrm{geo}}_{\mathcal{N}}
      \;=\; \varepsilon^{\mathrm{res}}_{\mathcal{N}}.
  \]
\end{lemma}
\begin{proof}
  \uses{mk-lem:modules_restrictFunctorIsoPullback_mathlib}
  The dual of \ref{mk-lem:gr_restrictAdjunction_unit_iso}: the counit of one adjunction
  is conjugated to the counit of the other by the uniqueness comparison of left
  adjoints.
\end{proof}

\begin{lemma}[Congruence cast at the pullback level]\label{mk-lem:gr_pullbackCongr_hom_eqToHom}
  \dcref{custom}

  For equal morphisms \(\varphi = \psi : X \to Y\) and a sheaf \(\mathcal{N}\) on \(Y\),
  the congruence isomorphism \(\varphi^{*}\mathcal{N} \cong \psi^{*}\mathcal{N}\)
  induced by \(\varphi = \psi\) is the canonical transport (the \(\mathrm{eqToHom}\) of
  the object-level equality \(\varphi^{*}\mathcal{N} = \psi^{*}\mathcal{N}\)).
\end{lemma}
\begin{proof}

Substituting \(\varphi = \psi\) reduces the congruence cast to the identity, which is
  the canonical transport of a reflexivity.
\end{proof}

\begin{lemma}[Trivial site-level congruence cast]\label{mk-lem:gr_restrictFunctorCongr_rfl}
  \dcref{custom}

  For an open immersion \(\varphi\) the site-level congruence cast associated to the
  reflexivity \(\varphi = \varphi\) is the identity natural transformation.
\end{lemma}
\begin{proof}

Pointwise on sections the cast is the presheaf restriction along a self-equality of
  opens, which is the identity by functoriality.
\end{proof}

\begin{lemma}[Comparison intertwines the two congruence casts]\label{mk-lem:gr_restrictIsoPullback_congr}
  \dcref{custom}

  \uses{mk-lem:gr_pullbackCongr_hom_eqToHom, mk-lem:gr_restrictFunctorCongr_rfl,
    mk-lem:modules_restrictFunctorIsoPullback_mathlib}
  For equal open immersions \(\varphi = \psi : X \to Y\) the left-adjoint comparison
  isomorphism intertwines the site-level and the pullback-level congruence casts:
  \[
    (\text{comparison}_\varphi)_{\mathcal{N}}\ \circ\
      (\varphi^{*}\!=\psi^{*})_{\mathcal{N}}
    \;=\;
    (\varphi^{\mathrm{res}}\!=\psi^{\mathrm{res}})_{\mathcal{N}}\ \circ\
      (\text{comparison}_\psi)_{\mathcal{N}}.
  \]
\end{lemma}
\begin{proof}
  \uses{mk-lem:gr_pullbackCongr_hom_eqToHom, mk-lem:gr_restrictFunctorCongr_rfl}
  Substituting \(\varphi = \psi\) collapses both congruence casts to identities
  (\ref{mk-lem:gr_pullbackCongr_hom_eqToHom}, \ref{mk-lem:gr_restrictFunctorCongr_rfl}),
  and the two comparison isomorphisms then coincide.
\end{proof}

\subsection{The candidate inverse}

\begin{definition}[Overlap base change in pullback form]\label{mk-def:gr_glueOverlapFactorIso}
  \dcref{custom}

  \uses{mk-lem:glueOverlapBaseChangeIso, mk-lem:modules_restrictFunctorIsoPullback_mathlib}
  Conjugating the overlap base change \(\beta_{ij}\)
  (\ref{mk-lem:glueOverlapBaseChangeIso}) on both sides by the comparison
  \ref{mk-lem:modules_restrictFunctorIsoPullback_mathlib} between the site-level
  restriction and the geometric pullback yields the isomorphism, in geometric form,
  \[
    \gamma_{ij} \;:\;
      \iota_i^{*}\bigl((\iota_j)_{*}\,\mathcal{M}_j\bigr)
        \;\xrightarrow{\ \sim\ }\;
      (f_{ij})_{*}\bigl((t_{ij}\circ f_{ji})^{*}\,\mathcal{M}_j\bigr),
  \]
  evaluated at the input sheaf \(\mathcal{M}_j\).
\end{definition}

\begin{definition}[The \(j\)-th component of the candidate inverse]\label{mk-def:gr_glueChartComponent}
  \dcref{custom}

  \uses{mk-def:gr_glueOverlapFactorIso}
  Given the transition isomorphisms \(g_{ij} : f_{ij}^{*}\mathcal{M}_i \xrightarrow{\sim}
  (t_{ij}\circ f_{ji})^{*}\mathcal{M}_j\), the \(j\)-th component of the candidate
  inverse \(\mathrm{s}_i\) is the composite
  \[
    \mathcal{M}_i
      \xrightarrow{\ \eta_{f_{ij}}\ }
    (f_{ij})_{*}\,f_{ij}^{*}\mathcal{M}_i
      \xrightarrow{\ (f_{ij})_{*}g_{ij}\ }
    (f_{ij})_{*}\,(t_{ij}\circ f_{ji})^{*}\mathcal{M}_j
      \xrightarrow{\ \gamma_{ij}^{-1}\ }
    \iota_i^{*}\bigl((\iota_j)_{*}\mathcal{M}_j\bigr),
  \]
  i.e. the geometric unit along \(f_{ij}\), the pushforward of the transition
  \(g_{ij}\), and the inverse of \(\gamma_{ij}\) (\ref{mk-def:gr_glueOverlapFactorIso}).
\end{definition}

\begin{definition}[The candidate-inverse family]\label{mk-def:gr_glueChartFamily}
  \dcref{custom}

  \uses{mk-def:gr_glueChartComponent, mk-def:gr_glue_equalizer}
  Assembling the components \ref{mk-def:gr_glueChartComponent} through the
  product-preservation comparison of \(\iota_i^{*}\) on the pushforward product
  \(P = \prod_j (\iota_j)_{*}\mathcal{M}_j\) (\ref{mk-def:gr_glue_equalizer}) gives a
  single morphism
  \[
    \mathcal{M}_i \;\longrightarrow\; \iota_i^{*}P,
  \]
  the candidate inverse before it is corrected to land in the equalizer
  \(\iota_i^{*}\mathcal{M}\).
\end{definition}

\begin{lemma}[The candidate family recovers each chart component]\label{mk-lem:gr_glueChartFamily_pullback_map_pi}
  \dcref{custom}

  \uses{mk-def:gr_glueChartComponent, mk-def:gr_glueChartFamily}
  For chart indices \(i,j\), composing the candidate family with the pullback of the
  \(j\)-th product projection recovers its \(j\)-th component:
  \[
    \mathrm{s}_i^{\mathrm{fam}}\circ
      \iota_i^*(\pi_j)=\mathrm{s}_i^{(j)}.
  \]
  This is the component formula for the product-preservation comparison used to
  define \ref{mk-def:gr_glueChartFamily}.
\end{lemma}

\subsection{The triple-overlap toolkit}

The first descent obligation (\ref{mk-lem:gr_glueChartFamily_equalizes}) reduces, factor
by factor of the restricted overlap product, to a single equation over a \emph{triple
overlap}. This subsection records that machinery. Fix three chart indices \(i, p, q\)
and form the triple overlap
\[
  V_{ipq} \;=\; V_{ip} \times_{U_i} V_{iq},
\]
with first projection \(\mathrm{fst} : V_{ipq} \to V_{ip}\). Two morphisms out of
\(V_{ipq}\) organise everything: the \emph{chart leg}
\(q_{pq} = \mathrm{fst}\circ f_{ip} : V_{ipq} \to U_i\) (down to \(U_i\)), and the
\emph{transport leg} \(\tau = t'_{ipq}\circ \mathrm{fst}_{pq} : V_{ipq} \to V_{pq}\)
(across to the pair overlap, through the cocycle transport \(t'_{ipq}\)). The square
\[
  \begin{array}{ccc}
    V_{ipq} & \xrightarrow{\ \ \tau\ \ } & V_{pq} \\
    \ \downarrow{\scriptstyle q_{pq}} & & \downarrow{\scriptstyle f_{pq}\circ \iota_p} \\
    U_i & \xrightarrow{\ \ \iota_i\ \ } & X
  \end{array}
\]
is the triple analogue of the chart--overlap square of \ref{mk-lem:glueOverlapBaseChangeIso}.

\begin{lemma}[Triple-overlap bridges of the glue datum]\label{mk-lem:gr_glueData_bridges}
  \dcref{custom}

  On the triple overlap \(V_{ijk} = V_{ij}\times_{U_i} V_{ik}\) the following three
  identities of morphisms hold, lining up the endpoints of the three base-change
  transports of the cocycle equation (C2):
  \begin{enumerate}
    \item \emph{(source)} the two projections of \(V_{ijk}\) to \(V_{ij}\) and \(V_{ik}\),
      followed by the overlap immersions \(f_{ij}\) and \(f_{ik}\) into \(U_i\),
      agree;
    \item \emph{(middle)} the \(ij\)-transport's target leg
      \((t_{ij}\circ f_{ji})\) and the \(jk\)-transport's source leg
      (through the cocycle transport \(t'_{ijk}\) and \(f_{jk}\)) agree as morphisms
      to \(U_j\);
    \item \emph{(target)} the composite \(jk\)-after-\(ij\) transport's target leg
      \((t_{jk}\circ f_{kj})\) and the \(ik\)-transport's target leg
      \((t_{ik}\circ f_{ki})\) agree as morphisms to \(U_k\).
  \end{enumerate}
\end{lemma}
\begin{proof}The source identity is the universal property (pullback condition) of \(V_{ijk}\).
  The middle identity follows from the transition factorisation
  \(t'_{ijk}\) in the glue datum together with the pullback condition.
  The target identity -- the heart of the cocycle -- combines the same transition
  factorisation, the pullback condition, the involutivity \(t_{ik}\circ t_{ki}=\mathrm{id}\)
  and the cocycle condition of the glue datum.
\end{proof}

\begin{lemma}[Pair component of the equalizing condition]\label{mk-lem:gr_glueChartComponent_leg_compat}
  \dcref{custom}

  \uses{mk-def:gr_glueChartComponent, mk-def:scheme_modules_glue,
    mk-lem:gr_glueData_bridges}
  For chart indices \(i,p,q\), the \(p\)-th candidate component followed by the
  restricted first descent-leg factor equals the \(q\)-th candidate component followed
  by the restricted second descent-leg factor.  After transposing to the triple overlap
  \(V_{ipq}\), this is precisely the cocycle multiplicativity condition (C2), with the
  endpoint identifications supplied by \ref{mk-lem:gr_glueData_bridges}.
\end{lemma}


\subsection{The first descent obligation}

\begin{lemma}[The candidate-inverse family equalizes the restricted legs (C2 transported)]\label{mk-lem:gr_glueChartFamily_equalizes}
  \dcref{custom}

  \uses{mk-def:gr_glueChartFamily, mk-def:scheme_modules_glue, mk-lem:gr_glueData_bridges,
    mk-lem:modules_pullback_basechange_transport, mk-lem:gr_overlap_appIso_compat,
    mk-lem:gr_glueChartComponent_self_counit, mk-lem:modules_restrictFunctorIsoPullback_mathlib}
  The family \ref{mk-def:gr_glueChartFamily} equalizes the two restricted descent legs
  \(\iota_i^{*}a\) and \(\iota_i^{*}b\); equivalently, it lands in the restricted
  equalizer \(\iota_i^{*}\mathcal{M}\). This is the first descent obligation, where the
  triple-overlap multiplicativity (C2) of \ref{mk-def:scheme_modules_glue} is consumed.
\end{lemma}
\begin{proof}
  \uses{mk-def:gr_glueChartFamily, mk-lem:gr_glueChartFamily_pullback_map_pi,
    mk-lem:gr_glueChartComponent_leg_compat}
    It suffices to check equality after projecting to each factor \((p,q)\) of the
  restricted overlap product, through the product-preservation comparison of
  \(\iota_i^{*}\) (the same comparison discipline as \ref{mk-lem:gr_glueRestrict_proj_compat}).
  Folding each restricted leg against the \((p,q)\)-projection -- splitting off the
  \(p\)- and \(q\)-th chart projections and cancelling the comparison via
  \ref{mk-lem:gr_glueChartFamily_pullback_map_pi} -- turns the \((p,q)\)-component into the
  equation
  \[
    \mathrm{s}_i^{(p)}\ \circ\ \iota_i^{*}\bigl(a_{(p,q)}\bigr)
      \;=\;
    \mathrm{s}_i^{(q)}\ \circ\ \iota_i^{*}\bigl(b_{(p,q)}\bigr),
  \]
  which is precisely the pair-component obligation
  \ref{mk-lem:gr_glueChartComponent_leg_compat}, where the triple-overlap multiplicativity
  (C2) is consumed. Assembling these componentwise equalities back through the comparison
  gives the equalizing condition.
\end{proof}

\begin{definition}[The candidate inverse \(\mathrm{s}_i\)]\label{mk-def:gr_glueRestrictionInv}
  \dcref{custom}

  \uses{mk-def:gr_glueChartFamily, mk-lem:gr_glueChartFamily_equalizes}
  Because \(\iota_i^{*}\) preserves the descent equalizer, the equalizing family
  \ref{mk-def:gr_glueChartFamily} (equalizing by \ref{mk-lem:gr_glueChartFamily_equalizes})
  lifts uniquely to a morphism
  \[
    \mathrm{s}_i \;:\; \mathcal{M}_i \;\longrightarrow\; \iota_i^{*}\mathcal{M},
  \]
  the candidate inverse of the chart restriction morphism.
\end{definition}

\begin{lemma}[Comparison compatibility of the restricted projection]\label{mk-lem:gr_glueRestrict_proj_compat}
  \dcref{custom}

  \uses{mk-def:gr_glue_equalizer, mk-def:scheme_modules_glue}
  Through the equalizer- and product-preservation comparisons of \(\iota_i^{*}\), the
  restricted equalizer inclusion of \(\iota_i^{*}\mathcal{M}\) into \(\iota_i^{*}P\)
  followed by the \(j\)-th product projection equals the chart pullback
  \(\iota_i^{*}(\pi_j)\) of the \(j\)-th descent-equalizer projection
  (\ref{mk-def:gr_glue_equalizer}).
\end{lemma}
\begin{proof}
  \uses{mk-def:gr_glue_equalizer}
  Pure limit bookkeeping: the equalizer- and product-comparison isomorphisms
  associated to a limit-preserving functor send the restricted inclusion and the
  restricted projection to the pullbacks of the inclusion and the projection
  respectively, and the composite of pullbacks is the pullback of the composite, which
  is \(\iota_i^{*}(\pi_j)\).
\end{proof}

\begin{lemma}[The candidate inverse recovers the components]\label{mk-lem:gr_glueRestrictionInv_proj}
  \dcref{custom}

  \uses{mk-def:gr_glueRestrictionInv, mk-lem:gr_glueRestrict_proj_compat,
    mk-def:gr_glueChartComponent}
  The candidate inverse \ref{mk-def:gr_glueRestrictionInv} followed by the restricted
  \(j\)-th projection \(\iota_i^{*}(\pi_j)\) recovers the \(j\)-th component
  \ref{mk-def:gr_glueChartComponent}:
  \[
    \mathrm{s}_i\ \circ\ \iota_i^{*}(\pi_j)
      \;=\; \mathrm{s}_i^{(j)} .
  \]
\end{lemma}
\begin{proof}
  \uses{mk-def:gr_glueRestrictionInv, mk-lem:gr_glueRestrict_proj_compat}
  Unfolding the lift \ref{mk-def:gr_glueRestrictionInv} and applying
  \ref{mk-lem:gr_glueRestrict_proj_compat}, the equalizer- and product-preservation
  comparisons cancel against the lift, leaving the defining \(j\)-th projection of the
  family \ref{mk-def:gr_glueChartFamily}, namely \(\mathrm{s}_i^{(j)}\).
\end{proof}

\begin{lemma}[Joint detection of morphisms into the restricted glued sheaf]\label{mk-lem:gr_glueRestrict_hom_ext}
  \dcref{custom}

  \uses{mk-lem:gr_glueRestrict_proj_compat}
  Two morphisms \(u, v : \mathcal{Z} \to \iota_i^{*}\mathcal{M}\) that agree after every
  restricted projection \(\iota_i^{*}(\pi_j)\) are equal.
\end{lemma}
\begin{proof}
  \uses{mk-lem:gr_glueRestrict_proj_compat}
  The restricted equalizer inclusion is a monomorphism and the restricted product
  projections are jointly monic, both through the preservation comparisons; so it
  suffices to test against \(\iota_i^{*}(\pi_j)\) for every \(j\), which is the
  hypothesis via \ref{mk-lem:gr_glueRestrict_proj_compat}.
\end{proof}

\begin{lemma}[Joint faithfulness of the chart restrictions]\label{mk-lem:gr_modules_glue_unique}
  \dcref{custom}

  \uses{mk-def:scheme_modules_glue}
  Two morphisms \(u, v : \mathcal{W} \to \mathcal{W}'\) of sheaves of modules on the glued
  scheme \(X\) of a glue datum (\ref{mk-def:scheme_modules_glue}) that agree after pullback
  \(\iota_i^{*}u = \iota_i^{*}v\) to every chart are equal. This is the characterisation of a
  glued sheaf up to unique isomorphism: a morphism out of the glued sheaf is determined by its
  chart-local restrictions.
\end{lemma}
\begin{proof}
  \uses{mk-def:scheme_modules_glue}
  This is the separation half of the sheaf condition over the chart cover \(\{\iota_i\}\):
  sections of \(\mathcal{W}'\) are detected on the chart-image opens \(\iota_i''(\iota_i^{-1}O)\),
  which cover any open \(O\) by the joint surjectivity of the chart immersions of the glue
  datum. Hence two morphisms agreeing on every chart agree on a covering family of opens, so
  they agree.
\end{proof}

\subsection{The three cocycle obligations}

\begin{lemma}[Self-component collapses to the identity (C1 + counit triangle)]\label{mk-lem:gr_glueChartComponent_self_counit}
  \dcref{custom}

  \uses{mk-def:gr_glueChartComponent, mk-def:scheme_modules_glue,
    mk-lem:gr_restrictIsoPullback_congr, mk-lem:gr_restrictAdjunction_unit_iso,
    mk-lem:gr_restrictIsoPullback_counit, mk-lem:gr_overlapBaseChange_inv_app,
    mk-lem:gr_pullbackCongr_hom_eqToHom}
  The \((i,i)\)-component \ref{mk-def:gr_glueChartComponent}, transposed back along the
  chart immersion \(\iota_i\) (post-composed with the geometric counit), is the
  identity of \(\mathcal{M}_i\):
  \[
    \mathrm{s}_i^{(i)}\ \circ\ \varepsilon^{\mathrm{geo}}_{\mathcal{M}_i}
      \;=\; \mathrm{id}_{\mathcal{M}_i}.
  \]
  This is the second descent obligation, consuming the self-identity (C1).
\end{lemma}
\begin{proof}
  \uses{mk-def:gr_glueChartComponent, mk-def:scheme_modules_glue,
    mk-lem:gr_restrictIsoPullback_congr, mk-lem:gr_restrictAdjunction_unit_iso,
    mk-lem:gr_restrictIsoPullback_counit, mk-lem:gr_overlapBaseChange_inv_app,
    mk-lem:gr_pullbackCongr_hom_eqToHom}
    Since \(f_{ii}\) factors as \(t_{ii}\circ f_{ii}\) (because \(t_{ii} = \mathrm{id}\)),
  (C1) of \ref{mk-def:scheme_modules_glue} says the transition \(g_{ii}\) is the
  canonical congruence cast; through \ref{mk-lem:gr_pullbackCongr_hom_eqToHom} and
  \ref{mk-lem:gr_restrictIsoPullback_congr} the three middle factors of the component
  collapse to a single site-level congruence cast. Rewriting the geometric unit by
  \ref{mk-lem:gr_restrictAdjunction_unit_iso} and the geometric counit by
  \ref{mk-lem:gr_restrictIsoPullback_counit} turns the whole composite into the
  site-level unit, the congruence cast, the section form of \(\beta_{ii}^{-1}\)
  (\ref{mk-lem:gr_overlapBaseChange_inv_app}), and the site-level counit. On sections
  over an open this is a cycle of presheaf restrictions along opens identities whose
  composite open-inclusion is the identity; folding the restrictions and using
  proof-irrelevance of the open inclusions gives the identity.
\end{proof}

\begin{lemma}[Transpose of the overlap base change (mate core)]\label{mk-lem:gr_glueOverlapFactor_transpose}
  \dcref{custom}

  \uses{mk-def:gr_glueOverlapFactorIso, mk-lem:glueOverlapBaseChangeIso,
    mk-lem:gr_overlapBaseChange_hom_app, mk-def:scheme_modules_glue}
  The adjoint transpose along \(\iota_i\) of the pullback-form overlap base change
  \(\gamma_{ij}\) (\ref{mk-def:gr_glueOverlapFactorIso}) is the canonical four-functor
  comparison of the glue square:
  \[
    \widehat{\gamma_{ij}}
      \;=\;
    (\iota_j)_{*}\bigl(\eta_{t_{ij}\circ f_{ji}}\bigr)
      \ \circ\ \mathrm{pushComp}
      \ \circ\ \mathrm{pushCongr}
      \ \circ\ \mathrm{pushComp}^{-1},
  \]
  the unit along \(t_{ij}\circ f_{ji}\) pushed forward along \(\iota_j\), regrouped by
  the pushforward-composition comparison, reindexed by the pushforward-congruence along
  the glue condition \((t_{ij}\circ f_{ji})\circ \iota_j = f_{ij}\circ \iota_i\), and
  ungrouped by the inverse pushforward-composition. The identity is free of the
  cocycle hypotheses: it is pure site-level content.
\end{lemma}
\begin{proof}
  \uses{mk-def:gr_glueOverlapFactorIso, mk-lem:glueOverlapBaseChangeIso,
    mk-lem:gr_overlapBaseChange_hom_app, mk-def:scheme_modules_glue,
    mk-lem:modules_restrictFunctorIsoPullback_mathlib}
    Expand the adjoint transpose of \(\gamma_{ij}\): the left-adjoint comparison
  isomorphisms (\ref{mk-lem:modules_restrictFunctorIsoPullback_mathlib}) conjugating the
  chart factor cancel against the unit/counit, leaving the bare site-level overlap base
  change \(\beta_{ij}\) (\ref{mk-lem:glueOverlapBaseChangeIso}). The unit bridge at
  \(t_{ij}\circ f_{ji}\) and the three naturality squares of the pushforward-composition
  and pushforward-congruence comparisons rewrite the transpose into the stated
  four-factor composite up to whiskering and reassociation. The residual site-level
  core is checked on sections over an open: by \ref{mk-lem:gr_overlapBaseChange_hom_app}
  the section form of \(\beta_{ij}\) is a presheaf restriction along an opens identity,
  and the remaining composite is a cycle of such restrictions whose total open
  inclusion is the identity, folded by proof-irrelevance exactly as in the closing step
  of \ref{mk-lem:gr_glueChartComponent_self_counit}.
\end{proof}

\begin{lemma}[Mate recognition of the factor transpose]\label{mk-lem:gr_glueOverlapFactor_mate}
  \dcref{custom}

  \uses{mk-lem:gr_glueOverlapFactor_transpose, mk-def:gr_glueOverlapFactorIso,
    mk-def:scheme_modules_glue}
  For any morphism \(m : \mathcal{W} \to (\iota_j)_{*}\mathcal{M}_j\) on the glued
  scheme, the conjugated transport of the inverse transpose of \(m\) along the glue
  square equals the inverse transpose, along \(f_{ij}\), of \(\iota_i^{*}(m)\) followed
  by \(\gamma_{ij}\) (\ref{mk-def:gr_glueOverlapFactorIso}). In other words, the factor
  transpose \ref{mk-lem:gr_glueOverlapFactor_transpose} is the mate of \(\gamma_{ij}\)
  through the pseudofunctor structure of pullback.
\end{lemma}
\begin{proof}
  \uses{mk-lem:gr_glueOverlapFactor_transpose, mk-def:gr_glueOverlapFactorIso,
    mk-def:scheme_modules_glue}
    Both sides are computed by transposing twice. By injectivity of the two transpose
  bijections it suffices to compare the double transposes. Naturality of the transpose
  in its left argument peels \(m\) off the right-hand side; the conjugate of the
  pullback-composition isomorphism along the glue condition is the inverse
  pushforward-composition isomorphism; and folding the unit pair and reindexing along
  the glue condition of \ref{mk-def:scheme_modules_glue} reduces the identity to the
  factor transpose \ref{mk-lem:gr_glueOverlapFactor_transpose}.
\end{proof}

\begin{lemma}[Overlap compatibility of the restriction morphisms]\label{mk-lem:gr_glueRestriction_overlap_compat}
  \dcref{custom}

  \uses{mk-lem:glueRestrictionHom, mk-def:scheme_modules_glue, mk-def:gr_glue_equalizer}
  The \((i,j)\)-component of the descent-equalizer condition, transposed to the
  pullback level, states that over the overlap \(V_{ij}\) the two restriction
  morphisms \(\mathrm{r}_i\) and \(\mathrm{r}_j\) (\ref{mk-lem:glueRestrictionHom}) differ
  exactly by the transition isomorphism \(g_{ij}\) (through the pseudofunctor
  congruence casts).
\end{lemma}
\begin{proof}
  \uses{mk-lem:glueRestrictionHom, mk-def:scheme_modules_glue, mk-def:gr_glue_equalizer}
  The glued sheaf is the descent equalizer, so its \(i\)- and \(j\)-projections satisfy
  the overlap-agreement condition of \ref{mk-def:gr_glue_equalizer}. Transposing that
  condition along \(\iota_i\) and \(\iota_j\) and folding the transposes back through
  the adjunction bijection turns it into the stated identity between the overlap
  pullbacks of \(\mathrm{r}_i\) and \(\mathrm{r}_j\), conjugated by the glue-condition
  casts of \ref{mk-def:scheme_modules_glue} and corrected by \(g_{ij}\).
\end{proof}

\begin{lemma}[Pair condition transposed: restriction composed with the component]\label{mk-lem:gr_glueRestrictionHom_glueChartComponent}
  \dcref{custom}

  \uses{mk-lem:gr_glueRestriction_overlap_compat, mk-lem:gr_glueOverlapFactor_mate,
    mk-def:gr_glueChartComponent, mk-lem:glueRestrictionHom}
  The chart restriction morphism \(\mathrm{r}_i\) (\ref{mk-lem:glueRestrictionHom})
  followed by the \(j\)-th candidate-inverse component \ref{mk-def:gr_glueChartComponent}
  equals the restricted \(j\)-th projection:
  \[
    \mathrm{r}_i\ \circ\ \mathrm{s}_i^{(j)}
      \;=\; \iota_i^{*}(\pi_j).
  \]
  This is the third descent obligation; together with
  \ref{mk-lem:gr_glueRestrictionInv_proj} and \ref{mk-lem:gr_glueRestrict_hom_ext} it gives
  \(\mathrm{r}_i \circ \mathrm{s}_i = \mathrm{id}\).
\end{lemma}
\begin{proof}
  \uses{mk-lem:gr_glueRestriction_overlap_compat, mk-lem:gr_glueOverlapFactor_mate,
    mk-def:gr_glueChartComponent, mk-lem:glueRestrictionHom}
    Unfold the component \ref{mk-def:gr_glueChartComponent}. By the overlap compatibility
  \ref{mk-lem:gr_glueRestriction_overlap_compat} the overlap pullback of \(\mathrm{r}_i\)
  transported across \(g_{ij}\) is the inverse mate transpose of the restricted
  projection \(\iota_i^{*}(\pi_j)\) through \(\gamma_{ij}\)
  (\ref{mk-lem:gr_glueOverlapFactor_mate}). Substituting this into the unit-naturality
  square of \(\mathrm{r}_i\) along \(f_{ij}\) and folding the resulting transpose
  cancels the factor isomorphism, leaving exactly \(\iota_i^{*}(\pi_j)\).
\end{proof}

\subsection{Joint faithfulness of chart restrictions}


\begin{definition}[The restriction isomorphism of the glued sheaf]\label{mk-def:glueRestrictionIso}
  \dcref{custom}

  \uses{mk-thm:isIso_glueRestrictionHom}
  Let \(\mathcal{M}\) be the glued sheaf of \ref{mk-def:scheme_modules_glue}, with the
  transition family \((g_{ij})\) satisfying (C1)/(C2). For each chart index \(i\) the
  restriction morphism \(\mathrm{r}_i\) (\ref{mk-lem:glueRestrictionHom}) is promoted to
  the \emph{restriction isomorphism}
  \[
    \rho_i \;:\; \iota_i^{*}\,\mathcal{M} \;\xrightarrow{\ \sim\ }\; \mathcal{M}_i,
  \]
  the canonical identification of the chart restriction of the glued sheaf with the
  \(i\)-th input sheaf, packaged from \ref{mk-thm:isIso_glueRestrictionHom}. Its
  underlying morphism is the adjoint transpose of the \(i\)-th descent-equalizer
  projection, and over each overlap \(V_{ij}\) the two restrictions \(\rho_i, \rho_j\)
  differ exactly by the transition isomorphism \(g_{ij}\).
\end{definition}
